% Generated draft: requires mathematical and source review.
\documentclass[11pt]{article}
\usepackage[margin=25mm]{geometry}
\usepackage{fontspec}
\setmainfont{Latin Modern Roman}
\usepackage{amsmath,amssymb,mathtools,unicode-math}
\setmathfont{Latin Modern Math}
\usepackage{enumitem,needspace,xurl,hyperref}
\setlength{\parindent}{0pt}
\setlength{\parskip}{5pt}
\setlength{\emergencystretch}{3em}
\newcommand{\N}{\mathbb{N}}
\newcommand{\Z}{\mathbb{Z}}
\newcommand{\Q}{\mathbb{Q}}
\newcommand{\R}{\mathbb{R}}
\newcommand{\C}{\mathbb{C}}
\newcommand{\F}{\mathbb{F}}
\newcommand{\A}{\mathbb{A}}
\newcommand{\PP}{\mathbb{P}}
\newcommand{\E}{\mathbb{E}}
\newcommand{\eps}{\varepsilon}
\newcommand{\dd}{\,\mathrm{d}}
\begin{document}
\section*{UnsolvedMath ID: 30002410}
\textbf{Source problem code:} catalogue entry
\subsection{Definitions and mathematical statement}
An involution constraint for an evolution PDE is understood here to mean a differential constraint on the unknown fields that is propagated by the evolution: initial data satisfying the constraint give rise to solutions satisfying it at subsequent times, whenever the requisite solution and propagation theory apply. Schematically, if the evolution is denoted by \(\mathcal E[u]=0\) and the constraint by \(\mathcal C[u]=0\), propagation means
\[
\mathcal C[u](\cdot,0)=0
\quad\Longrightarrow\quad
\mathcal C[u](\cdot,t)=0
\]
throughout the time interval of existence, subject to appropriate regularity and compatible boundary conditions. This notation explains the terminology only; the supplied record does not specify \(\mathcal E\), \(\mathcal C\), a spatial domain, a solution class, or initial and boundary conditions.

The source asks the following research questions for involution-constrained PDEs arising in magnetohydrodynamics and nonlinear elasticity:
\begin{enumerate}
\item How can the proper involutions of the PDE system be identified?
\item What is the relation between enforcing these involutions in a numerical method and using methods of higher order of accuracy?
\item Are methods of lower order of accuracy sufficient when involution constraints are enforced?
\end{enumerate}
Here, numerical enforcement means requiring an appropriate discrete counterpart of a constraint to hold, with the discrete constraint and the meaning of enforcement still to be specified. The terms ``proper,'' ``higher-order,'' ``lower-order,'' and ``sufficient'' are not given mathematical criteria in the record. In particular, no accuracy order, error norm, convergence requirement, stability property, or application-specific performance target is prescribed. The source also does not specify whether enforcement is exact or approximate.

Consequently, these are structural and constructive research questions rather than a fully quantified conjecture. A precise theorem-level formulation would require the PDE systems, admissible constraints, analytical setting, numerical methods, and criterion of sufficiency to be supplied. No universal necessity or sufficiency claim is asserted here.
\subsection{Short English statement}
The problem asks how to identify the relevant differential constraints in magnetohydrodynamics and nonlinear elasticity, how enforcing them interacts with numerical accuracy, and whether constraint enforcement can make lower-order methods adequate. The supplied record leaves the systems, numerical setting, and meaning of adequacy unspecified, so the draft preserves the questions without asserting a precise equivalence or sufficiency theorem.
\subsection{Sources}
\begin{itemize}
\item Supplied local UnsolvedMath registry record, permanent ID \(30002410\), title ``Involutions and Accuracy in Constrained PDE Methods,'' source problem code \texttt{catalogue entry}, in the collection ``Oberwolfach Reports Open Problems.'' The record gives proposed year \(2013\) and names no individual proposer.
\item Supplied report citation: \emph{High-Resolution Mathematical and Numerical Analysis of Involution-Constrained PDEs} (2014), Oberwolfach Reports, pp.~2691--2747, DOI \texttt{10.4171/owr/2013/47}. Supplied source link: \url{https://doi.org/10.4171/owr/2013/47}. The background describes these reports as documenting mathematical workshops and their open-problem discussions.
\item Supplied literature source, ``Oberwolfach Report'': \url{https://ems.press/content/serial-article-files/46478}.
\item Supplied literature source, ``Dumbser--Chiocchetti--Peshkov'': \url{https://arxiv.org/abs/2003.02799}.
\item Supplied literature source, ``Balsara--Kumar--Chandrashekar'': \url{https://doi.org/10.2140/camcos.2021.16.1.59}.
\item Literature assessment dated \texttt{2026-08-21}: the supplied curation labels the problem open and states that numerical methods exist for particular curl-involution systems, while the extracted broad structural and constructive questions remain unsettled. This is supplied dated curation derived from the OWR material, not an independently verified literature assessment. Its time-sensitive status and cited claims require independent verification.
\item The registry labels the statement ``exact'' and reports manual review finding the extracted statement self-contained, mathematically coherent, and free of LaTeX or report-layout artifacts. These are supplied source-verification claims; they do not resolve the missing mathematical specifications identified in this draft.
\end{itemize}
\end{document}
